\chapter*{Abstract}

Simultaneous interpretation is a cornerstone of multilingual communication in international conferences and institutional settings, yet it traditionally relies on highly skilled human interpreters and specialized hardware infrastructures that are costly, resource-intensive, and difficult to scale. Recent advances in automatic speech recognition (ASR), neural machine translation (NMT), and text-to-speech (TTS) synthesis have made AI-based interpretation technically feasible; however, achieving reliable real-time performance remains a challenging engineering problem due to strict latency, stability, and integration constraints. This thesis investigates the design, implementation, and evaluation of a real-time AI-based simultaneous interpretation system from a system-engineering perspective. The proposed solution adopts a modular, provider-agnostic architecture that integrates streaming ASR, NMT, and TTS synthesis within a unified pipeline, with particular emphasis on incremental processing, stability-aware segmentation, asynchronous execution, and multi-language output support. The objectives are threefold: (i) to design and validate a real-time interpretation pipeline capable of processing continuous speech input with bounded end-to-end latency, (ii) to develop a system-level evaluation framework that quantifies latency and stability trade-offs, (iii) and to assess the feasibility of integrating AI-generated interpretation output into professional interpretation hardware environments. Experimental results obtained under standardized real-time conditions demonstrate that the proposed system maintains latency within acceptable bounds for simultaneous interpretation across multiple configurations, while component-level analysis highlights recognition and segmentation as the dominant contributors to delay. Ablation studies reveal non-trivial trade-offs between responsiveness and output stability, underscoring the importance of system-level orchestration. Finally, a hardware integration feasibility study shows that AI-based interpretation output can be aligned with professional channel-based audio infrastructures using standardized digital audio formats, supporting practical deployment beyond software-only demonstrations.